\section{Limitaciones y trabajo futuro}

Limitaciones asumidas en esta iteración:
\begin{itemize}
  \item la plataforma modela un subconjunto activo de DCAT-AP-ES con caso HVD, no el perfil completo con todos los recomendados,
  \item la conformidad estricta sin \texttt{Warning} no es objetivo actual; se priorizan obligatorios y consistencia reproducible,
  \item CKAN no se usa como origen operativo porque cosechar un catálogo externo replica metadatos ya publicados y no evidencia suficientemente gobierno desde sistemas fuente,
  \item la publicación externa del catálogo no se automatiza en esta iteración,
  \item sin requisitos enterprise (HA, SSO, hardening, observabilidad avanzada).
\end{itemize}

Trabajo futuro propuesto:
\begin{itemize}
  \item ampliar algunos metadatos recomendados para reducir advertencias SHACL,
  \item orquestar el workflow con Airflow u otro planificador para exportar y, si procede, subir periódicamente el catálogo a un CKAN externo,
  \item enviar por correo a responsables de catálogo el informe de validación RDF/SHACL,
  \item generar informes presentables en PDF o HTML además de JSON y TTL,
  \item exportar también RDF/Turtle además de JSON-LD,
  \item estudiar CKAN como destino de publicación o federación, no como fuente principal de generación de metadatos.
\end{itemize}
